\documentclass[a4paper,12pt]{article}

\usepackage[a4paper,margin=2cm]{geometry}
\usepackage[T1]{fontenc}
\usepackage[utf8]{inputenc}
\usepackage[ngerman,english]{babel}
\usepackage{xcolor}
\usepackage[most]{tcolorbox}
\usepackage{array}
\usepackage{parskip}
\usepackage{booktabs}
\usepackage{colortbl}

\definecolor{HeaderBlueDark}{RGB}{70,110,170}
\definecolor{RowBlueLight}{RGB}{235,243,255}
\definecolor{RowBlueDark}{RGB}{220,233,250}

\setlength{\parindent}{0pt}
\linespread{1.15}

\begin{document}
    
    \section{Die Präposition und das Wort \glqq zu\grqq}
        
        \textbf{zu} ist ein sehr wichtiges Wort im Deutschen. Es kann als Präposition, als Teil einer Infinitivkonstruktion und auch als Adverb benutzt werden. Die wichtigste Bedeutung ist oft: \textbf{Richtung zu einer Person, einem Ort oder einem Ziel}. Als Präposition steht \textbf{zu} fast immer mit dem \textbf{Dativ}.
        
        \begin{tcolorbox}[colback=RowBlueLight,colframe=HeaderBlueDark,title=\textbf{Grundregel}]
            \textbf{zu + Dativ = to; toward; to the place of}\\
            ex: Ich gehe zu Klaus: I am going to Klaus's place.\\
            ex: Ich gehe zum Arzt: I am going to the doctor.\\
            ex: Ich gehe zur Arbeit: I am going to work.
        \end{tcolorbox}
    
    
    \section{Formen: zum, zur, zu den}
        
        \textbf{zu} verbindet sich oft mit dem Artikel. Diese Formen sind sehr wichtig:
        
        \rowcolors{2}{RowBlueLight}{RowBlueDark}
        \begin{tabular}{>{\bfseries}p{4cm}p{5cm}p{5cm}}
            \rowcolor{HeaderBlueDark}
            \color{white}\textbf{Form} & \color{white}\textbf{Deutsch} & \color{white}\textbf{English} \\
            zu + dem                   & zum Arzt                      & to the doctor                 \\
            zu + der                   & zur Schule                    & to school                     \\
            zu + den                   & zu den Kindern                & to the children               \\
            zu + einem                 & zu einem Freund               & to a friend                   \\
            zu + einer                 & zu einer Freundin             & to a female friend            \\
        \end{tabular}
    
    
    \section{zu bei Personen}
        
        Wenn man zu einer Person geht, benutzt man meistens \textbf{zu}. Sehr oft bedeutet das: zu dieser Person nach Hause oder an den Ort, wo diese Person ist.
        
        ex: Ich gehe zu Klaus: I am going to Klaus's place.\\
        ex: Kommst du heute zu mir?: Are you coming to my place today?\\
        ex: Wir fahren morgen zu meinen Eltern: We are going to my parents tomorrow.\\
        ex: Sie geht zu ihrer Freundin: She is going to her friend's place.
        
        \textbf{Wichtig:} Wenn man schon dort ist, benutzt man oft \textbf{bei}.
        
        ex: Ich bin bei Klaus: I am at Klaus's place.\\
        ex: Ich war gestern bei meinen Eltern: I was at my parents' place yesterday.
    
    
    \section{zu bei Orten, Terminen und Institutionen}
        
        Man benutzt \textbf{zu} oft, wenn man zu einem Ziel, zu einem Termin oder zu einer Institution geht. Dabei ist nicht immer wichtig, ob man wirklich in ein Gebäude hineingeht. Wichtig ist das Ziel.
        
        ex: Ich gehe zum Arzt: I am going to the doctor.\\
        ex: Ich muss zur Bank: I have to go to the bank.\\
        ex: Wir fahren zum Bahnhof: We are going to the train station.\\
        ex: Er geht zur Arbeit: He is going to work.\\
        ex: Ich gehe zum Unterricht: I am going to class.\\
        ex: Sie geht zur Party: She is going to the party.
    
    
    \section{zu, nach und in: der Unterschied}
        
        Viele Deutschlernende verwechseln \textbf{zu}, \textbf{nach} und \textbf{in}. Die einfache Regel ist:
        
        \rowcolors{2}{RowBlueLight}{RowBlueDark}
        \begin{tabular}{>{\bfseries}p{3cm}p{5.5cm}p{5.5cm}}
            \rowcolor{HeaderBlueDark}
            \color{white}\textbf{Wort} & \color{white}\textbf{Benutzung}                  & \color{white}\textbf{Beispiel} \\
            zu                         & Personen, Termine, Institutionen, Ziel allgemein & Ich gehe zum Arzt.             \\
            nach                       & Städte, Länder ohne Artikel, Hause               & Ich fahre nach Wien.           \\
            in                         & Bewegung in einen Raum oder in ein Gebiet        & Ich gehe ins Haus.             \\
        \end{tabular}
        
        ex: Ich fahre nach Österreich: I am going to Austria.\\
        ex: Ich fahre nach Wien: I am going to Vienna.\\
        ex: Ich gehe nach Hause: I am going home.\\
        ex: Ich gehe zu Klaus nach Hause: I am going to Klaus's home.\\
        ex: Ich gehe ins Haus: I am going into the house.
        
        \textbf{Nicht richtig:} Ich gehe nach Klaus Hause.\\
        \textbf{Richtig:} Ich gehe zu Klaus.\\
        \textbf{Richtig:} Ich gehe zu Klaus nach Hause.
    
    
    \section{zu + Infinitiv}
        
        \textbf{zu} steht auch vor dem Infinitiv eines Verbs. Dann ist \textbf{zu} keine normale Präposition, sondern ein Teil der Infinitivkonstruktion.
        
        ex: Ich versuche, Deutsch zu lernen: I try to learn German.\\
        ex: Es ist wichtig, pünktlich zu sein: It is important to be punctual.\\
        ex: Ich habe vergessen, dich anzurufen: I forgot to call you.\\
        ex: Sie hat keine Zeit, einkaufen zu gehen: She has no time to go shopping.
        
        Bei trennbaren Verben steht \textbf{zu} zwischen Präfix und Verb:
        
        ex: anrufen $\rightarrow$ anzurufen\\
        ex: einkaufen $\rightarrow$ einzukaufen\\
        ex: aufstehen $\rightarrow$ aufzustehen\\
        ex: anfangen $\rightarrow$ anzufangen
        
        ex: Ich versuche, früher aufzustehen: I try to get up earlier.\\
        ex: Er hat vergessen, mich anzurufen: He forgot to call me.
    
    
    \section{um ... zu, ohne ... zu, statt ... zu}
        
        Es gibt auch feste Konstruktionen mit \textbf{zu}.
        
        \textbf{um ... zu} bedeutet: in order to. Man benutzt es für ein Ziel oder einen Zweck.
        
        ex: Ich lerne Deutsch, um in Österreich zu arbeiten: I learn German in order to work in Austria.\\
        ex: Ich spare Geld, um eine Kamera zu kaufen: I save money in order to buy a camera.
        
        \textbf{ohne ... zu} bedeutet: without doing something.
        
        ex: Er ging weg, ohne etwas zu sagen: He left without saying anything.\\
        ex: Ich habe gearbeitet, ohne eine Pause zu machen: I worked without taking a break.
        
        \textbf{statt ... zu} oder \textbf{anstatt ... zu} bedeutet: instead of doing something.
        
        ex: Ich bleibe zu Hause, statt auszugehen: I stay at home instead of going out.\\
        ex: Sie schreibt eine Nachricht, anstatt anzurufen: She writes a message instead of calling.
    
    
    \section{zu als Adverb: too}
        
        \textbf{zu} kann auch \textbf{too} bedeuten. Dann steht es vor einem Adjektiv oder Adverb und zeigt, dass etwas mehr ist, als gut oder passend ist.
        
        ex: Das Auto ist zu teuer: The car is too expensive.\\
        ex: Der Kaffee ist zu heiß: The coffee is too hot.\\
        ex: Du fährst zu schnell: You are driving too fast.\\
        ex: Ich habe zu wenig Zeit: I have too little time.\\
        ex: Es gibt zu viele Probleme: There are too many problems.
    
    
    \section{Feste Ausdrücke mit zu}
        
        Es gibt viele feste Ausdrücke mit \textbf{zu}. Diese sollte man als ganze Ausdrücke lernen.
        
        \rowcolors{2}{RowBlueLight}{RowBlueDark}
        \begin{tabular}{>{\bfseries}p{4cm}p{5cm}p{5cm}}
            \rowcolor{HeaderBlueDark}
            \color{white}\textbf{Ausdruck} & \color{white}\textbf{Bedeutung} & \color{white}\textbf{Beispiel}  \\
            zu Fuß                         & on foot                         & Ich gehe zu Fuß.                \\
            zum Beispiel                   & for example                     & Zum Beispiel lerne ich Deutsch. \\
            zu Hause                       & at home                         & Ich bin zu Hause.               \\
            nach Hause                     & home, toward home               & Ich gehe nach Hause.            \\
            zu spät                        & too late                        & Ich komme zu spät.              \\
            zu Ende                        & finished, over                  & Der Film ist zu Ende.           \\
            zu Besuch                      & visiting                        & Ich bin bei Klaus zu Besuch.    \\
            zu zweit                       & as two people                   & Wir arbeiten zu zweit.          \\
            bis zu                         & up to                           & Das kostet bis zu 100 Euro.     \\
        \end{tabular}
    
    
    \section{Kurze Zusammenfassung}
        
        \textbf{zu} hat mehrere wichtige Funktionen. Als Präposition steht es mit dem Dativ und bedeutet oft \textbf{to} oder \textbf{toward}. Man benutzt es besonders bei Personen, Terminen, Institutionen und Zielen.
        
        ex: Ich gehe zu Klaus: I am going to Klaus's place.\\
        ex: Ich gehe zum Arzt: I am going to the doctor.\\
        ex: Ich fahre zur Arbeit: I am going to work.
        
        Mit Infinitiv bedeutet \textbf{zu} ungefähr \textbf{to} vor einem Verb.
        
        ex: Ich versuche, Deutsch zu lernen: I try to learn German.
        
        Als Adverb bedeutet \textbf{zu} oft \textbf{too}.
        
        ex: Der Kaffee ist zu heiß: The coffee is too hot.
        
        Die wichtigste Regel ist: \textbf{zu + Dativ}. Außerdem muss man den Unterschied zwischen \textbf{zu}, \textbf{nach} und \textbf{in} gut lernen, weil diese Wörter im Deutschen nicht genau wie im Englischen benutzt werden.

\end{document}